\section{Probability extras: characteristic functions, modes, one-sided BM}
\label{sec:prob}

These are not red-syllabus holes. They are the process and limit objects that 2025--2026 never asked in this form, and that close the remaining official headings ``characteristic function'', ``modes of convergence'', and ``basic properties of Brownian motion / random walks''.

\subsection{Characteristic-function CLT with infinite variance}

\begin{definition}[Characteristic function]
\label{def:cf}
\(\phi_{X}(t)=\mathbb{E}[e^{\mathrm{i}tX}]\). It determines the law, is uniformly continuous, and converts independence into products. L\'evy's continuity theorem: \(\phi_{n}(t)\to\phi(t)\) pointwise with \(\phi\) continuous at \(0\) if and only if \(X_{n}\Rightarrow X\) with \(\phi_{X}=\phi\).
\end{definition}

The classical Lindeberg CLT (already used as a truncation argument in 2026 Spring P2) assumes a finite variance. 2024 Spring P1 is the CF sibling: \(\mathbb{E}X^{2}=\infty\), a given small-\(t\) expansion \(1-\phi(t)\sim t^{2}\lvert\log\lvert t\rvert\rvert\), and a sequence \(a_{n}\) such that \(S_{n}/a_{n}\Rightarrow N(0,1)\). The computation is
\[
\phi_{S_{n}/a_{n}}(t)=\phi(t/a_{n})^{n}
=
\exp\bigl(n\log\phi(t/a_{n})\bigr)
\sim
\exp\bigl(-t^{2}/2\bigr)
\]
once \(a_{n}\) is chosen so that \(n(1-\phi(t/a_{n}))\to t^{2}/2\). For the displayed expansion one takes \(a_{n}\sim\sqrt{n\log n}\).

Do not answer this with Lindeberg. The hypothesis \(\mathbb{E}X^{2}=\infty\) is the point.

\subsection{Modes of convergence}

Almost-sure, in probability, in \(L^{p}\), and in distribution are already used as language. The remaining named fact is the algebra of the modes.

\begin{proposition}
\label{prop:modes}
If \(X_{n}\to X\) and \(Y_{n}\to Y\) in probability, then \(X_{n}Y_{n}\to XY\) in probability. The same statement is false in \(L^{1}\): take \(X_{n}=Y_{n}=n\mathbf{1}_{[0,1/n]}\) on \([0,1]\), or any pair whose product has expectation one while each factor tends to zero in \(L^{1}\).
\end{proposition}

This is 2024 Spring P3. Connection: a.s.\ convergence is also a ring (products and continuous maps); \(L^{1}\) is a Banach space but not an algebra under pointwise products. Convergence in distribution is not preserved by products unless one factor is tight and converges in probability (Slutsky).

\subsection{One-sided Brownian hitting and biased walks}

2025 Spring P5 is two-sided hitting \(T_{a}=\inf\{t:\lvert B_{t}\rvert=a\}\) and independence of \(\operatorname{sign}(B_{T_{a}})\). The missing transform is the one-sided time \(T_{a}=\inf\{t:B_{t}=a\}\) for \(a>0\).

The exponential martingale \(M_{t}^{u}=\exp(uB_{t}-u^{2}t/2)\) is bounded up to \(T_{a}\) if \(u>0\). Optional stopping and \(B_{T_{a}}=a\) give
\[
\mathbb{E}[e^{-u^{2}T_{a}/2}]=e^{-ua},
\qquad
\mathbb{E}[e^{-\lambda T_{a}}]=e^{-a\sqrt{2\lambda}}.
\]
That is 2024 Spring P6. The same martingale, with a drift, produces the ruin probability of a biased random walk (2024 Spring P4): if \(\mathbb{E}[e^{\theta\xi}]=1\) for a nonzero \(\theta\), then \(e^{\theta S_{n}}\) is a martingale and
\[
\mathbb{P}(T_{k}<\infty)=e^{-\theta k}
\]
on the side against the drift. \(\sup_{n}S_{n}\) is then geometric.

2023 Fall P2 is the discrete sibling: on a Laplace density \(f(x)=\tfrac12 e^{-\lvert x\rvert}\) one builds \(M_{n}(\theta)\) by the same exponential tilt. The cumulant is
\[
A_{n}(\theta)
=
n\log\mathbb{E}[e^{\theta X_{1}}]
=
n\log\Bigl(\frac{1}{1-\theta^{2}}\Bigr),
\qquad
\lvert\theta\rvert<1,
\]
and then \(M_{n}(\theta)=\exp(\theta S_{n}-A_{n}(\theta))\) is a martingale. 2025--2026 used exponential martingales for geometric Brownian motion and for weighted walks; they did not ask you to \emph{construct} the tilt from a density.

\subsection{Exponential tails and exponential extremes}

\paragraph{Tails \(\Leftrightarrow\) moments \(\Leftrightarrow\) MGF.}
\paper{2024 Fall P3.} For a real random variable the following are equivalent up to absolute constants \(K_{i}\): an exponential tail \(\mathbb{P}(\lvert X\rvert\ge t)\le 2e^{-t/K_{1}}\); the moment bound \(\mathbb{E}\lvert X\rvert^{p}\le(K_{2}p)^{p}\); and \(\mathbb{E}\exp(\lvert X\rvert/K_{3})\le 2\). The script is integration by tails one way and Stirling \(n!\asymp(n/e)^{n}\) the other way. Sub-Gaussian tails are the same machine with \(t^{2}\) in the exponent; do not quote a CLT.

\paragraph{Exponential maxima.}
\paper{2024 Spring P2.} Let \(X_{k}\) be i.i.d.\ unit exponential. Borel--Cantelli on the independent events \(\{X_{n}>(1\pm\varepsilon)\log n\}\) gives \(\limsup X_{n}/\log n=1\) almost surely, and therefore \(M_{n}/\log n\to 1\) almost surely for \(M_{n}=\max_{k\le n}X_{k}\). The 2025--2026 papers never asked an extreme-value one-liner.

\subsection{Connections and how examined}

CF problems give a small-\(t\) expansion and ask for \(a_{n}\). Mode problems ask for a proof and a counter-example, not a slogan. BM problems ask you to name the martingale, justify stopping, and read the Laplace transform. Biased-walk problems ask for the root of the MGF, the ruin probability, and the law of the supremum.

Possible questions are exactly the three anchors below, plus ``construct \(M_{n}(\theta)\) on a given density so that it is a martingale''.

{\footnotesize
\begin{center}
\begin{tabular}{@{}L{2.7cm}L{5.4cm}L{6.1cm}@{}}
\toprule
Anchor & What the paper wants & Near-miss \\
\midrule
\gap{2024 Spring P1} & \(\mathbb{E}X^{2}=\infty\); CF \(\Rightarrow\) \(S_{n}/a_{n}\Rightarrow N(0,1)\) & 2026 Spring P2 is finite-variance truncation \\
\gap{2024 Spring P3} & products in probability; \(L^{1}\) counter-example & language of a.s.\ / in distribution \\
\gap{2024 Spring P6} & Laplace of one-sided \(T_{a}\) via \(e^{uB-u^{2}t/2}\) & 2025 Spring P5 is two-sided + sign \\
\gap{2024 Spring P4} & biased walk: exp.\ martingale, ruin, \(\sup S_{n}\) & symmetric hitting already done \\
\gap{2023 Fall P2} & build \(M_{n}(\theta)\) on a Laplace density & GBM exponential martingale \\
\gap{2024 Fall P3} & exp.\ tail \(\Leftrightarrow\) moments \(\Leftrightarrow\) MGF & finite-variance CLT \\
\gap{2024 Spring P2} & \(X_{n}/\log n\) and \(M_{n}/\log n\) & order statistics of uniforms \\
\bottomrule
\end{tabular}
\end{center}
}
